\newpage
\section*{Capitulo 7. Análisis del sistema}
\addcontentsline{toc}{section}{Capitulo 7. Análisis del sistema}
\setcounter{section}{7}
\setcounter{subsection}{0}

\subsection{Requerimientos del sistema}
El presente proyecto tiene como producto principal el desarrollo de un sistema móvil para la detección de insectos descortezadores mediante visión por computadora. Sin embargo, durante el análisis del problema se identificó que una aplicación móvil por sí sola no es suficiente para cubrir todas las necesidades del proceso de monitoreo forestal: mientras que el experto de campo requiere una herramienta ágil, funcional sin conexión a internet y orientada a la captura y análisis inmediato en sitio, los gestores y analistas del bosque necesitan acceder al historial completo de registros, consultar estadísticas, revisar los resultados oficiales del modelo en la nube y generar reportes de seguimiento.
Por esta razón, el sistema se compone de dos subsistemas complementarios: una aplicación móvil Android diseñada bajo un enfoque offline-first, que constituye el núcleo del prototipo, y una plataforma web que actúa como herramienta de monitoreo y gestión para el personal de administración del bosque. Ambos subsistemas comparten la misma base de datos en la nube y están diseñados para operar de forma coordinada.

\subsubsection{Requerimientos funcionales}
Los requisitos funcionales definen las capacidades y comportamientos específicos que el sistema debe satisfacer, diferenciados por componente: la aplicación móvil, orientada al monitoreo en campo, y el panel web, destinado a la gestión y consulta de resultados. 

Con el propósito de reflejar claramente a qué subsistema aplica cada requerimiento, se ha adoptado un esquema de etiquetado de requisitos funcionales con el ámbito de aplicación correspondiente: aplicación móvil (\textbf{MOV}), aplicación web (\textbf{WEB}) o ambos entornos (\textbf{MOV/WEB}).

La Tabla 3 resume estos requisitos.

\begin{longtable}{|>
{\centering\arraybackslash}m{1.8cm}|>{\centering\arraybackslash}m{2.2cm}|m{10cm}|}
\caption{Requerimientos funcionales de la aplicación móvil y la plataforma web}\label{tab:rf} \\
\hline
\textbf{ID} & \textbf{Ámbito} & \textbf{Descripción} \\
\hline
\endfirsthead
\multicolumn{3}{c}%
{{\bfseries Tabla \thetable\ -- continuación}} \\
\hline
\textbf{ID} & \textbf{Ámbito} & \textbf{Descripción} \\
\hline
\endhead

RF-01 & MOV/WEB & El sistema debe permitir al usuario registrarse e iniciar sesión utilizando un correo electrónico y una contraseña. \\
\hline
RF-02 & MOV/WEB & El sistema debe permitir al usuario solicitar la recuperación de su contraseña a través de un enlace enviado a su correo electrónico. \\
\hline
RF-03 & MOV & El sistema debe permitir la entrada de imágenes de dos formas: capturándolas directamente desde la aplicación o subiéndolas desde la galería del dispositivo. \\
\hline
RF-04 & MOV & El sistema debe contar con herramientas de edición básicas antes de enviar la imagen a procesar. \\
\hline
RF-05 & MOV & El sistema debe registrar automáticamente la ubicación geográfica del árbol fotografiado, obtenida a partir de los metadatos de la fotografía o del GPS del dispositivo en el momento de la captura. \\
\hline
RF-06 & MOV & El sistema debe almacenar localmente un registro de todas las imágenes ingresadas por el usuario, incluyendo su resultado de clasificación, \textit{timestamp} y ubicación, para su posterior sincronización con la nube. \\
\hline
RF-07 & MOV & El sistema debe sincronizar los registros almacenados localmente hacia la nube cuando el usuario cuente con conexión a internet, conservando la fecha y hora original de captura como referencia del registro. \\
\hline
RF-08 & MOV/WEB & La aplicación móvil deberá incorporar un modelo en formato TensorFlow Lite embebido en el dispositivo para realizar una clasificación local e inmediata de la imagen capturada cuando no se tenga acceso a internet.
La aplicación web deberá incorporar un modelo en formato Keras3 para realizar una clasificación en la nube de la imagen capturada en la aplicación movil.  
\\
\hline
RF-09 & MOV & El sistema debe emitir y mostrar en pantalla una clasificación inmediata del estado sanitario del árbol, indicando si está sano o infectado, generada por el modelo local embebido en el dispositivo, sin requerir conexión a internet. \\
\hline
RF-10 & WEB & El sistema debe permitir agrupar y dar seguimiento a un conjunto de muestras específicas a lo largo del tiempo, generando un historial del estado de salud de cada ejemplar. \\
\hline
RF-11 & MOV/WEB & El sistema debe mostrar alertas o indicadores visuales con el analisis complementario de color que faciliten la detección de focos de infestación. \\
\hline
RF-12 & MOV & El sistema debe mostrar un resumen estadístico de la jornada actual, incluyendo el número total de fotografías tomadas en el día, la cantidad de árboles clasificados como sanos y la cantidad de posibles infecciones detectadas. \\
\hline
RF-13 & MOV & El sistema debe proveer acceso directo y visible a la función principal: ``Nueva Captura''. \\
\hline






RF-14 & MOV & El sistema debe mostrar notificaciones o alertas cuando el usuario se encuentre en las proximidades de una zona previamente catalogada como foco de infestación. \\
\hline







RF-15 & MOV & La sincronización manual estará disponible únicamente cuando existan errores de sincronización y el dispositivo cuente con conexión activa. En caso contrario, el sistema deberá deshabilitar esta opción e informar al usuario la causa. \\
\hline
RF-16 & WEB & El sistema debe permitir generar un reporte básico en formato PDF de las capturas realizadas. \\
\hline
RF-17 & WEB & El sistema debe permitir descargar el reporte generado. \\
\hline
RF-18 & MOV & El sistema debe conservar en el almacenamiento local del dispositivo los registros asociados a las imágenes capturadas, excepto estas últimas. \\
\hline
RF-19 & MOV & El sistema debe permitir al usuario asociar una nueva captura a un árbol ya existente en el registro local, o bien crear un nuevo registro de árbol al momento de realizar una captura. \\
\hline
RF-20 & WEB & El sistema debe mostrar el historial completo y permanente de todos los árboles registrados, incluyendo sus capturas, resultados de clasificación, observaciones y fechas de subida y sincronización, sin restricción de periodo de tiempo. \\
\hline






RF-21 & WEB & El sistema debe mostrar estadísticas generales del estado del arbolado registrado, incluyendo totales por estado sanitario, distribución por especie y evolución temporal de los focos de infestación detectados. \\
\hline


RF-22 & MOV/WEB & El sistema movil debe ejecutar un análisis complementario de color sobre la imagen capturada y mostrar su resultado como información complementaria al resultado de clasificación, tanto en la app móvil como en la app web. \\
\hline
RF-23 & MOV & El sistema debe permitir al usuario operativo editar la información de un registro de árbol, incluyendo la actualización de su estado sanitario y la adición de observaciones o notas de campo. \\
\hline
RF-24 & WEB & El sistema debe mostrar un mapa general con la ubicación de todos los árboles registrados, permitiendo filtrarlos por estado sanitario, especie, fecha de captura y dispositivo. \\
\hline

RF-25 & MOV & El sistema debe mostrar un mapa general con la ubicación de todos los árboles registrados del dispositivo, permitiendo filtrarlos por estado sanitario, especie y fecha de captura. \\
\hline

RF-26 & WEB & El sistema debe permitir al administrador gestionar las cuentas de usuario del sistema, incluyendo la creación, edición, desactivación y asignación de roles. \\
\hline

RF-27 & MOV/WEB & El sistema deberá permitir que una cuenta de experto de campo pueda estar asociada a uno o varios dispositivos móviles. Cada dispositivo deberá mantener una base de datos local independiente en SQLite para almacenar árboles, capturas, análisis y ubicaciones registrados offline. Al recuperar conexión, la aplicación móvil deberá sincronizar los registros locales con la base de datos central en la nube, conservando la identificación del dispositivo de origen. \\
\hline

RF-28 & MOV & Cuando dos dispositivos asociados a la misma cuenta registren árboles con el mismo nombre, el sistema no deberá asumir que se trata del mismo árbol físico. \\
\hline

\end{longtable}


\subsubsection{Requerimientos no funcionales}
Los requisitos no funcionales establecen las restricciones de calidad, rendimiento y operación bajo las cuales el sistema debe funcionar, independientemente de las funcionalidades que provee. La Tabla 4 los detalla.

\begin{longtable}{|>{\centering\arraybackslash}m{1.8cm}|>{\centering\arraybackslash}m{2.5cm}|m{9.5cm}|}
\caption{Requerimientos no funcionales de la aplicación móvil y la plataforma web}\label{tab:rnf} \\
\hline
\textbf{ID} & \textbf{Categoría} & \textbf{Descripción} \\
\hline
\endfirsthead

\multicolumn{3}{c}%
{{\bfseries Tabla \thetable\ -- continuación}} \\
\hline
\textbf{ID} & \textbf{Categoría} & \textbf{Descripción} \\
\hline
\endhead

RNF-01 & Rendimiento & El sistema debe devolver el resultado del análisis (sano/infectado) en un tiempo máximo de X segundos una vez que la imagen ha sido cargada exitosamente en la nube. \\
\hline
RNF-02 & Portabilidad & El tamaño de la aplicación móvil instalada no debe exceder los 50 MB, asegurando que pueda ser instalada en dispositivos de gama media sin saturar el almacenamiento. \\
\hline
RNF-03 & Usabilidad & Las transiciones entre las pantallas de la aplicación, como pasar del mapa al historial del árbol, no deben demorar más de X segundos. \\
\hline
RNF-04 & Eficiencia & La aplicación debe estar optimizada para no realizar procesos en segundo plano innecesarios, minimizando el consumo de batería del dispositivo durante las jornadas de inspección en el bosque. \\
\hline
RNF-05 & Seguridad & Toda comunicación entre la aplicación móvil y los servicios en la nube, incluyendo el envío de credenciales, imágenes y coordenadas, debe realizarse a través de protocolos seguros como HTTPS. \\
\hline
RNF-06 & Seguridad & Las contraseñas de los usuarios deben almacenarse en la base de datos utilizando algoritmos de hashing para evitar su exposición. \\
\hline
RNF-07 & Seguridad & El sistema debe cerrar automáticamente la sesión del usuario tras un periodo prolongado de inactividad, por ejemplo 24 horas, para proteger los datos en caso de extravío del dispositivo. \\
\hline
RNF-08 & Precisión & El modelo de visión por computadora debe alcanzar y mantener un porcentaje mínimo de precisión de X\% en la correcta clasificación de los árboles infestados por descortezadores. \\
\hline
RNF-09 & Arquitectura & El sistema debe respetar una arquitectura cliente-servidor, delegando el procesamiento pesado del modelo de inteligencia artificial a la nube. \\
\hline
RNF-10 & Persistencia & El sistema debe contar con una base de datos relacional o no relacional optimizada para manejar consultas rápidas del historial forestal y referenciar los archivos de imágenes. \\
\hline
RNF-11 & Implementación & La aplicación cliente debe ser desarrollada de forma nativa para el sistema operativo Android utilizando el lenguaje de programación \textbf{Kotlin}, garantizando un acceso óptimo al hardware del dispositivo, como cámara, GPS y almacenamiento local. \\
\hline
RNF-12 & Implementación & El modelo de visión por computadora para la detección de insectos descortezadores debe ser programado en Python, utilizando \textit{frameworks} especializados en aprendizaje profundo, como TensorFlow o PyTorch. \\
\hline
RNF-13 & Rendimiento & Debido a que el análisis de imágenes puede tomar algunos segundos, el backend en Python debe manejar las peticiones de forma asíncrona para no bloquear la interfaz de usuario en la aplicación Kotlin. \\
\hline
RNF-14 & Mantenibilidad & El proyecto móvil debe utilizar Gradle para la gestión de librerías, mientras que el entorno de Python debe usar gestores como \textit{pip} o \textit{conda} para asegurar la replicabilidad del entorno de desarrollo. \\
\hline
RNF-15 & Persistencia & La aplicación cliente en Android debe implementar \textbf{SQLite}, preferentemente mediante la librería \textit{Room}, para el almacenamiento estructurado y local de la información. Esto permitirá guardar temporalmente los datos de las capturas cuando el usuario se encuentre en zonas sin conexión a internet, garantizando el uso básico \textit{offline}. \\
\hline
RNF-16 & Persistencia & El backend del sistema debe utilizar \textbf{PostgreSQL} como motor de base de datos relacional principal. Esta base de datos será la encargada de almacenar de forma centralizada, segura y escalable las credenciales de acceso de los usuarios, así como el historial completo y digitalizado del estado sanitario de cada árbol evaluado. \\
\hline
RNF-17 & Integridad & El sistema debe cambiar el estado local de los registros en SQLite a ``Sincronizado'' o eliminarlos localmente una vez que el servidor backend confirme que la información se guardó exitosamente en PostgreSQL, evitando así la duplicidad de datos. \\
\hline
\end{longtable}

\subsection{Reglas de negocio del sistema}
Las reglas de negocio definen los criterios que rigen el comportamiento del sistema en relación con la detección de plagas y el procesamiento de imágenes. Reflejan decisiones metodológicas independientes de la tecnología utilizada y se presentan en la Tabla 5.

\begin{longtable}{|>{\centering\arraybackslash}m{1.8cm}|>{\centering\arraybackslash}m{2.5cm}|m{9.5cm}|}
\caption{Reglas de negocio de la aplicación móvil y la plataforma web}\label{tab:reglas_operacion} \\
\hline
\textbf{ID} & \textbf{Tipo} & \textbf{Descripción} \\
\hline
\endfirsthead

\multicolumn{3}{c}%
{{\bfseries Tabla \thetable\ -- continuación}} \\
\hline
\textbf{ID} & \textbf{Tipo} & \textbf{Descripción} \\
\hline
\endhead

RN-01 & Acceso & Solo los usuarios con rol \textbf{Administrador} podrán crear cuentas nuevas. El registro no será abierto al público. \\
\hline
RN-02 & Acceso & El rol \textbf{Experto de campo} podrá iniciar sesión en la app móvil y web. El rol \textbf{Administrador} solo podrá iniciar sesión en la app web. \\
\hline
RN-03 & Sesión & La sesión en la app móvil persistirá activa hasta que el usuario cierre sesión manualmente. No expirará por inactividad. \\
\hline
RN-04 & Seguridad & El enlace de recuperación de contraseña tendrá una vigencia máxima de 24 horas a partir de su generación. \\
\hline
RN-05 & Seguridad & El enlace de recuperación será de un solo uso; una vez utilizado quedará invalidado automáticamente. \\
\hline
RN-06 & Seguridad & El sistema permitirá un máximo de 3 solicitudes de recuperación de contraseña por usuario en un periodo de 30 días naturales. \\
\hline
RN-07 & Imagen & Solo se aceptarán imágenes en formato JPG, JPEG, PNG, WEBP, BMP y HEIC. \\
\hline
RN-08 & Imagen & El tamaño máximo de imagen aceptado será de 10 MB. \\
\hline
RN-09 & Imagen & La resolución mínima aceptada será de 300$\times$300 píxeles, correspondiente al tamaño de entrada requerido por EfficientNet-B3. Imágenes con resolución menor serán rechazadas por el sistema con un mensaje informativo al usuario. \\
\hline
RN-10 & Edición & La herramienta de edición disponible antes del análisis será recorte de la imagen. \\
\hline
RN-11 & Edición & La herramienta de edición se aplicará únicamente antes del análisis. Una vez procesada la imagen, no podrá modificarse. \\
\hline
RN-12 & Edición & La edición de la imagen es opcional; el usuario podrá omitirla y proceder directamente al análisis. \\
\hline
RN-13 & Ubicación & Si la imagen proviene de la galería y contiene metadatos de ubicación, el sistema utilizará esas coordenadas como ubicación del árbol. \\
\hline
RN-14 & Ubicación & Si la imagen no contiene metadatos de ubicación, el sistema utilizará las coordenadas del GPS del dispositivo en el momento del registro. \\
\hline
RN-15 & Ubicación & Si no es posible obtener la ubicación por ninguna de las dos vías, el registro se guardará sin coordenadas y se marcará como \textbf{Ubicación no disponible}, sin bloquear el flujo de captura. \\
\hline
RN-16 & Registros & El usuario podrá eliminar manualmente un registro local antes de sincronizarlo. Esta acción será irreversible y requerirá confirmación explícita. \\
\hline
RN-17 & Registros & Cada registro local tendrá uno de los siguientes estados: \textbf{Pendiente de sincronización}, \textbf{Sincronizado} o \textbf{Error de sincronización}. \\
\hline
RN-18 & Registros & Cada registro local deberá contener como mínimo: imagen, resultado del modelo, timestamp de captura, coordenadas GPS y árbol asociado. \\
\hline
RN-19 & Sincronización & La sincronización automática se activará cuando el dispositivo detecte conexión a internet, ya sea por datos móviles o Wi-Fi. \\
\hline
RN-20 & Sincronización & En caso de fallo durante la sincronización, el sistema reintentará automáticamente hasta 3 veces antes de marcar el registro como \textbf{Error de sincronización}. \\
\hline
RN-21 & Sincronización & La fecha y hora de sincronización se almacenará como un campo adicional separado del timestamp original de captura. Ambos valores serán visibles en la app web. \\
\hline
RN-22 & Modelo & El resultado generado por el modelo se mostrará al experto inmediatamente después de la captura, etiquetado explícitamente como \textbf{Offline}. \\
\hline

RN-23 & Modelo & El sistema clasificará el árbol en dos categorías únicamente: \textbf{Sano} o \textbf{Infestado}. \\
\hline

% Sujeto a cambios en %
RN-24 & Modelo & El modelo TFLite solo emitirá una clasificación si el nivel de confianza supera el 75\%. Si no supera ese umbral, el resultado se mostrará como \textbf{No determinado} y se recomendará al experto retomar la captura. \\
\hline
RN-25 & Alertas & De acuerdo con la NOM-019-SEMARNAT-2017 \cite{nom019_semarnat_2017}, se considerará \textbf{foco de infestación activo} cuando se registren 3 o más árboles clasificados como infectados dentro de una superficie de 1{,}000 m$^2$. \\
\hline
RN-26 & Interfaz & Los indicadores visuales de estado sanitario utilizarán el siguiente código de color: \textbf{verde} para sano, \textbf{rojo} para infestado y \textbf{amarillo} para no determinado. \\
\hline







RN-27 & Alertas & De acuerdo con la NOM-019-SEMARNAT-2017 \cite{nom019_semarnat_2017}, el sistema considerará que el experto está en proximidad de un foco de infestación cuando se encuentre a 100 metros o menos de las coordenadas de un árbol catalogado como foco. \\
\hline
RN-28 & Alertas & La alerta de proximidad se generará usando únicamente el GPS del dispositivo y los registros locales almacenados en SQLite, funcionando sin conexión a internet. \\
\hline
RN-29 & Alertas & No se emitirá más de una alerta por foco en un periodo de 60 minutos para evitar notificaciones repetitivas mientras el experto permanezca en la zona. \\













\hline
RN-30 & Jornada & El resumen de jornada considerará únicamente los registros cuyo timestamp de captura corresponda al día en curso, desde las 00:00 hasta las 23:59. \\
\hline
RN-31 & Jornada & El resumen de jornada se actualizará en tiempo real conforme el experto realice nuevas capturas durante la jornada. \\
\hline
RN-32 & Sincronización & El sistema debe permitir al usuario reintentar manualmente la sincronización de los registros marcados con estado \textbf{Error de sincronización}, siempre que el dispositivo cuente con conexión a internet activa. \\
\hline
RN-33 & Reportes & Los reportes generados desde la app web incluirán el historial sincronizado en la nube de un periodo de tiempo o de una zona específica. \\
\hline
RN-34 & Reportes & Todo reporte deberá incluir como mínimo: identificador del árbol, nombre del árbol, fecha y hora de captura, resultado de clasificación, coordenadas GPS y nombre del experto que realizó la captura. \\
\hline
RN-35 & Almacenamiento & La eliminación automática de registros del almacenamiento local solo podrá aplicarse a registros con estado \textbf{Sincronizado}. Los registros en estado \textbf{Pendiente} o \textbf{Error} no deberán eliminarse automáticamente hasta que sean sincronizados o eliminados manualmente por el usuario. \\
\hline






RN-36 & Historial & El historial visible en la app móvil mostrará únicamente los registros del dia actual. \\
\hline








RN-37 & Árboles & No puede existir una captura sin un árbol asociado. La selección o creación del árbol es un paso obligatorio antes de proceder al análisis. \\
\hline
RN-38 & Árboles & Al crear un nuevo árbol, el sistema generará automáticamente un identificador único. El usuario deberá asignar manualmente un nombre al árbol y seleccionar su especie de una lista predefinida basada en las especies hospederas documentadas para el Bosque de Chapultepec, entre ellas \textbf{Fresno} (\textit{Fraxinus uhdei}) y \textbf{Cedro blanco} (\textit{Cupressus lusitanica}), además de otras que el Administrador podrá gestionar desde la app web. \\
\hline
RN-39 & Árboles & El nombre asignado al árbol podrá complementarse con una descripción libre opcional, por ejemplo: \textit{Pino 1 con bordes irregulares en corteza}, para facilitar su identificación en campo. \\
\hline
RN-40 & Usuarios & Solo el rol \textbf{Administrador} podrá crear, editar, desactivar y asignar roles a usuarios. \\
\hline
RN-41 & Usuarios & El \textbf{Experto de campo} podrá editar registros de árboles y agregar observaciones en la app movil. \\
\hline
RN-42 & Historial & El historial en la app web no tendrá restricción de tiempo; mostrará todos los registros sincronizados desde el inicio del sistema. \\
\hline









RN-43 & Alertas & El radio de proximidad para alertas definido en RN-29 será configurable únicamente desde la app web por el \textbf{Administrador}. \\
\hline










RN-44 & Alertas & De acuerdo con la NOM-019-SEMARNAT-2017 \cite{nom019_semarnat_2017}, el sistema deberá considerar como periodos de mayor riesgo los meses de enero a mayo y de septiembre a noviembre, pudiendo emitir alertas estacionales en la app web durante estos periodos. \\
\hline










RN-45 & Modelo & Las evidencias visuales de infestación que el sistema buscará identificar mediante el modelo de clasificación incluyen: presencia de grumos de resina rojiza en el fuste, cambios de coloración del follaje de verde a verde amarillento o rojizo, y presencia visible de galerías o daños en la corteza. \\
\hline








% VER CUAL MLIB ES
RN-46 & Análisis cromático & El análisis complementario de color se ejecutará de forma local en el dispositivo de forma concurrente al modelo TensorFlow Lite, utilizando OpenCV, sin requerir conexión a internet. \\
\hline
RN-47 & Análisis cromático & El análisis de color del follaje clasificará los píxeles en tres rangos cromáticos correspondientes a las fases de infestación documentadas: verde, verde-amarillento o verde limón, y café-rojizo. El resultado se expresará como porcentaje de cada rango sobre el total de píxeles analizados. \\
\hline
RN-48 & Análisis cromático & El análisis de color de la corteza clasificará los píxeles en los rangos cromáticos asociados a las fases de colonización subcortical: blanco/crema/rojizo, café-anaranjado y café-grisáceo. El resultado se expresará como porcentaje de cada rango sobre el total de píxeles analizados. \\












\hline
RN-49 & Análisis cromático & El análisis cromático se ejecutará de forma independiente según el tipo de imagen procesada, ya sea follaje o corteza, dado que cada nivel opera en un espacio visual distinto. No se realizarán comparaciones ni agregaciones cromáticas entre imágenes de diferente naturaleza dentro de un mismo análisis. \\
\hline

RN-50 & Almacenamiento & La base de datos local SQLite de la aplicación móvil no deberá almacenar imágenes como datos binarios de forma permanente. Únicamente deberá guardar metadatos del análisis y una referencia, ruta temporal o URL de la evidencia fotográfica.\\
\hline


RN-51 & Historial & La evidencia fotográfica completa solo estará disponible para consulta desde la plataforma web después de que el registro haya sido sincronizado correctamente con la nube.\\
\hline

RN-52 & Sincronización & La sincronización de datos se realizará únicamente desde la base local del dispositivo móvil hacia la base central en la nube. El sistema no contemplará la descarga automática de registros desde la nube hacia otros dispositivos móviles.\\
\hline

RN-53 & Historial & La aplicación web permitirá consultar y analizar los registros sincronizados desde los dispositivos móviles, pero no modificará los registros operativos generados en la app móvil.\\
\hline

RN-54 & Usuarios & Las funciones CRUD del Administrador aplicarán únicamente sobre datos administrativos o de configuración del sistema almacenados en la nube.\\
\hline

\end{longtable}

\subsection{Catálogo de mensajes del sistema}
La Tabla~\ref{tab:catalogoMensajes} presenta el catálogo de mensajes utilizados por el sistema para informar al usuario sobre errores, validaciones o estados relevantes durante la ejecución de los casos de uso.

\begin{longtable}{|p{2.5cm}|p{10.5cm}|}
\caption{Catálogo de mensajes del sistema.}
\label{tab:catalogoMensajes} \\

\hline
\rowcolor{black}
\textcolor{white}{\textbf{Código}} & \textcolor{white}{\textbf{Mensaje}} \\ \hline
\endfirsthead

\hline
\rowcolor{black}
\textcolor{white}{\textbf{Código}} & \textcolor{white}{\textbf{Mensaje}} \\ \hline
\endhead

\hline
\endfoot

\hline
\endlastfoot

MSG1 & Dato requerido faltante. \\ \hline
MSG2 & El correo electrónico no está registrado. \\ \hline
MSG3 & Contraseña incorrecta. \\ \hline
MSG4 & La cuenta se encuentra desactivada. \\ \hline
MSG5 & El dispositivo no está autorizado para esta cuenta. \\ \hline
MSG6 & No fue posible enviar el correo de recuperación. \\ \hline
MSG7 & El enlace o código de recuperación no es válido. \\ \hline
MSG8 & El enlace o código de recuperación ha expirado. \\ \hline
MSG9 & Las contraseñas no coinciden. \\ \hline
MSG10 & La contraseña no cumple con los requisitos de seguridad. \\ \hline
MSG11 & No existe una sesión activa. \\ \hline
MSG12 & No fue posible cerrar la sesión correctamente. \\ \hline
MSG13 & No existen registros para la jornada actual. \\ \hline
MSG14 & No fue posible cargar la información local. \\ \hline
MSG15 & No existen registros con el estado sanitario seleccionado. \\ \hline
MSG16 & No fue posible cargar los registros recientes. \\ \hline
MSG17 & No fue posible acceder a la cámara. \\ \hline
MSG18 & No fue posible importar la imagen seleccionada. \\ \hline
MSG19 & Imagen no válida para el análisis. \\ \hline
MSG20 & No fue posible registrar la ubicación GPS. \\ \hline
MSG21 & No fue posible guardar el análisis en la base local. \\ \hline
MSG22 & No fue posible obtener la clasificación del estado sanitario. \\ \hline
MSG23 & No fue posible recortar la imagen. \\ \hline
MSG24 & Área de recorte no válida. \\ \hline
MSG25 & No fue posible sincronizar el registro. \\ \hline
MSG26 & No existen registros sincronizados disponibles. \\ \hline
MSG27 & No fue posible cargar la información consolidada. \\ \hline
MSG28 & No fue posible aplicar los filtros seleccionados. \\ \hline
MSG29 & No se cuenta con permisos de administración. \\ \hline
MSG30 & No fue posible cargar la información de usuarios. \\ \hline
MSG31 & El correo electrónico ya se encuentra registrado. \\ \hline
MSG32 & No fue posible guardar los cambios de la cuenta. \\ \hline
MSG33 & El dispositivo ya se encuentra asociado a una cuenta. \\ \hline
MSG34 & No fue posible asociar el dispositivo. \\ \hline
MSG35 & No fue posible desasociar el dispositivo. \\ \hline

\end{longtable}


\subsection{Actores del sistema}
Un actor representa cualquier entidad externa que interactúa con el sistema, ya sea una persona, un rol organizacional o un sistema externo. En el contexto de UML, los actores no forman parte del sistema en sí, sino que se relacionan con él a través de los casos de uso, iniciando acciones o recibiendo resultados de las mismas.

Para el presente sistema, los actores fueron definidos considerando los roles operativos involucrados en el proceso de monitoreo forestal, así como los sistemas externos con los que ambos subsistemas se comunican.

Por consiguiente, se identificaron los siguientes actores:

\begin{enumerate}
    \item \textbf{Experto de campo}\\
    Persona encargada de utilizar la aplicación móvil en campo y la plataforma web para consultar la información sincronizada. Este actor registra árboles, realiza análisis del estado sanitario y consulta los datos generados durante la jornada.
        \begin{itemize}
            \item Iniciar sesión en la aplicación móvil o web.
            \item Recuperar y restablecer contraseña.
            \item Capturar o importar imágenes de árboles.
            \item Recortar o repetir una imagen antes del análisis.
            \item Asociar una imagen a un árbol nuevo o existente.
            \item Registrar datos de árboles nuevos.
            \item Modificar datos de árboles existentes desde la aplicación móvil.
            \item Obtener la clasificación del estado sanitario del árbol.
            \item Consultar el análisis complementario de color.
            \item Consultar registros recientes de la jornada.
            \item Consultar el historial local de análisis de la jornada.
            \item Consultar el mapa local de árboles.
            \item Filtrar registros, árboles o mapas según el estado sanitario u otros criterios disponibles.
            \item Consultar notificaciones y alertas.
            \item Reintentar manualmente la sincronización de registros con error.
            \item Consultar información consolidada desde la plataforma web.
            \item Consultar historial general, mapa general, estadísticas y perfil de usuario en la web.
            \item Cerrar sesión.
        \end{itemize}

    \item \textbf{Administrador}\\
    Persona encargada de gestionar las cuentas de usuario y los dispositivos asociados desde la plataforma web. Este actor no modifica los registros operativos capturados desde la aplicación móvil, sino que administra la información de acceso y configuración relacionada con las cuentas.
        \begin{itemize}
            \item Iniciar sesión en la plataforma web.
            \item Recuperar y restablecer contraseña.
            \item Crear cuentas de usuario.
            \item Editar cuentas de usuario.
            \item Desactivar cuentas de usuario.
            \item Asignar roles de usuario.
            \item Gestionar dispositivos asociados a una cuenta.
            \item Asociar dispositivos móviles a una cuenta.
            \item Desasociar dispositivos móviles de una cuenta.
            \item Consultar dispositivos asociados a una cuenta.
            \item Cerrar sesión.
        \end{itemize}

    \item \textbf{Sistema}\\
    Componente encargado de ejecutar procesos automáticos internos necesarios para el funcionamiento de la aplicación móvil y la plataforma web.
        \begin{itemize}
            \item Almacenar registros en la base local SQLite del dispositivo móvil.
            \item Sincronizar registros locales hacia la base de datos en la nube cuando exista conexión a internet.
            \item Registrar el estado de sincronización de cada análisis.
            \item Marcar registros como sincronizados tras confirmación del servidor.
            \item Marcar registros con error de sincronización cuando corresponda.
            \item Consolidar en la nube los registros enviados desde los dispositivos asociados a una cuenta.
            \item Generar estadísticas, mapas y vistas consolidadas para la plataforma web.
            \item Identificar focos de infestación activos a partir de los registros sincronizados.
            \item Generar notificaciones de sincronización y alertas sanitarias.
        \end{itemize}
\end{enumerate}



\subsection{Métricas de evaluación del proyecto}
\subsubsection{Métricas de desempeño del modelo}
\subsubsection{Métricas de desempeño del sistema}
\subsubsection{Criterios de aceptación del prototipo}

\subsection{Análisis de riesgos y restricciones}
\subsubsection{Riesgos asociados al conjunto de datos}
\subsubsection{Riesgos técnicos de implementación}
\subsubsection{Restricciones del prototipo}

\newpage
\subsubsection{Medidas de mitigación}

\subsection{Plan general de pruebas}
\subsubsection{Pruebas del modelo de clasificación}
\subsubsection{Pruebas del sistema}
\subsubsection{Pruebas de integración}
